\begin{cvsection}{Master Thesis}
    \cvitem
    {Discreteness Effects in N-body Simulations of Cosmological Structure Formation}
    [Mar 2017]
    [\ipa]
    [Supervisor: Prof.\ Dr.\ Cristiano Porciani]
    [Cosmological N-body simulations solve the collisionless Boltzmann equation in an approximate way by sampling phase-space with discrete particles and computing the gravitational interactions between them. This (macro-)particles are many orders of magnitude more massive than any plausible dark-matter candidate and it is thus necessary to fine tune the force resolution in order to suppress few-body relaxation effects. All this might corrupt the output with unphysical components ranging from the contamination of physical quantities with Poisson noise to the creation of spurious halos. In this Master\'s thesis, we quantify the importance of discreteness effects by finding and comparing ``matching-halos'' in a set of cosmological simulations with different mass and force resolutions but identical initial conditions. Our results show a tight relationship between the number of particles in a virialized halo and the accuracy of the most of its properties.]
\end{cvsection}